\chapter{TÍNH TOÁN, THIẾT KẾ VÀ CHẾ TẠO HỆ THỐNG}
\label{chap:tombo-L1}

This chapter focuses on the application of the $\mathcal{L}_1$Quad controller to the Tombo quadrotor, enabling effective handling of mild and short-duration collisions occurring while the drone follows a predefined trajectory towards a target point. The $\mathcal{L}_1$Quad architecture, grounded in geometric control combined with an $\mathcal{L}_1$ adaptive control strategy, compensates for uncertainties such as collision-induced disturbances that cause trajectory tracking errors at each time step. This compensation is achieved through a piecewise constant adaptation law, which allows fast and reliable error correction during flight. Experimental results demonstrate that the integration of the Tombo quadrotor with $\mathcal{L}_1$Quad effectively mitigates collision impacts, especially those occurring at the propeller, by rapidly compensating tracking errors and thereby reducing collision recovery time.

\section{Quadrotor Specifications}\label{sec:section1}
\subsection{Quadrotor Hardware}
A custom-built quadrotor platform, shown in Figure~\ref{fig: tombo quadrotor l1}, is constructed using an FPV QAV-RXL 2 – 10" frame (Lumenier). The propulsion system consists of four SURPASS 2216 1120KV \gls{BLDC} (Goolsky), each individually driven by a 30A \gls{ESC} (RC Electric Parts). Power is supplied by a 1500 mAh 4-cell (4S) 95C \gls{LiPo} battery (Lumenier), which is distributed to the ESCs for simultaneous motor control. The \gls{FCU} is a Pixhawk 6X (Holybro), which handles low-level stabilization and communicates with the onboard controller for trajectory tracking and collision handling. The total mass of the quadrotor, including all components, is approximately 1.18\,kg. The distance between the two front motors is $l$=0.345\,m, while the distance between the front and rear motors is $d$=0.26\,m.
\begin{figure}[H]
    \centering
    \includegraphics[scale=0.3]{Figures/tombo quadrotor l1.pdf}
    \captionsetup{width=\textwidth,justification=centering}
    \caption{The Tombo quadrotor.}
    \label{fig: tombo quadrotor l1}
\end{figure}

The evaluation of the proposed $\mathcal{L}_1$Quad controller in collision scenarios is conducted using two types of propellers: (1) commercially available rigid propellers and (2) Tombo propellers. These propellers are interchangeable on the same quadrotor platform to ensure consistency in experimental comparison. The rigid propellers used for comparison are 9x3.8SF-B4 (APC), which features a length of 9 inches and a pitch of 3.8 inches. To ensure a fair evaluation, Tombo propellers are designed to match the same length as rigid propellers, while differing in structural design to allow for deformation upon impact. In particular, the rigid propeller is thinner and lighter than the Tombo propeller, weighing 9\,g compared to 18\,g. This mass difference reflects the added material and structure in the Tombo design to enable its deformability.

\subsection{Measurement of the moment of inertia of quadrotor}
To experimentally determine the moment of inertia of the quadrotor along each principal axis, the bifilar pendulum method was employed. In this approach, the drone is suspended by two parallel strings of equal length attached symmetrically about the axis of interest. As shown in Figure~\ref{fig: moment of inertia measurement}, the experimental setup is designed to measure the moment of inertia $J_{zz}$ about the rotational axis $b_3$ in the body frame of the quadrotor. When the drone is rotated slightly about this axis and released, it undergoes torsional oscillations. By recording the oscillation period and knowing the mass of the drone, the distance between the suspension strings, and the length of the strings, the moment of inertia $J$ can be calculated using the following relation:
\begin{equation}
    J = \frac{m g D^2 T^2}{4 \pi^2 L},
\end{equation}
where $m$ is the mass of the drone, $g$ is the gravitational acceleration, $D$ is the distance from the quadrotor's center of mass to each suspension string, $L$ is the length of the strings, and $T$ is the oscillation period. This procedure was repeated for each of the three principal axes to obtain the corresponding inertia values $J_{xx}$, $J_{yy}$, and $J_{zz}$ as shown in Table~\ref{fig: inertia moment}.
\begin{figure}[ht]
      \centering
      \includegraphics[scale=0.5]{Figures/drone inertia moment experimental setup.pdf}
      \caption{Experimental setup of the bifilar pendulum for $J_{zz}$ moment of inertia estimation.}
      \label{fig: moment of inertia measurement}
\end{figure}
\begin{table}[H]
  \centering 
  \setlength{\tabcolsep}{0pt}
  \caption[Measured quadrotor inertia moments]{Experimentally measured moments of inertia of the quadrotor with respect to the body-fixed rotational axes.}
  \begin{tabular}{|C{3.2cm}|C{3.5cm}|C{2.8cm}|}
    \hline
    Axis & Inertia moment & Value (kgm$^2$) \\
    \hline
    $b_1$ & $J_{xx}$ &  0.0126 \\
    \hline
    $b_2$ & $J_{yy}$ &  0.0081 \\
    \hline
    $b_3$ & $J_{zz}$ &  0.0181 \\
    \hline
  \end{tabular}
  \label{fig: inertia moment}
\end{table}

\subsection{Thrust and moment measurement of propeller and motor}
In quadrotor control, regulating the thrust generated by the four propellers is essential to achieve and maintain stable flight. Each propeller produces thrust and generates torque, directly affecting the attitude and yaw control of the vehicle. Therefore, precise control requires accurate modeling of both thrust and torque contributions from all four motors. Since the thrust and torque produced by each motor-propeller pair depend on the \gls{PWM} duty cycle supplied through the \gls{ESC}, it is crucial to experimentally determine the relationship between \gls{PWM} input and the resulting forces and moments. This characterization enables more accurate force allocation and attitude control, which are fundamental for achieving stability, responsiveness, and robustness in flight performance.
\begin{figure}[ht]
      \centering
      \includegraphics[scale=0.5]{Figures/Thrust force measurement setup.pdf}
      \caption[Thrust measurement setup]{Experimental setup for measuring propeller thrust force and moment as a function of \gls{PWM} percentage from the Ardupilot controller.}
      \label{fig: thrust force measurement setup}
\end{figure}

An experimental setup, illustrated in Figure~\ref{fig: thrust force measurement setup}, was designed to determine the relationship between the \gls{PWM} percentage supplied to the motor and the resulting thrust and moment, based on the specific motor and propeller hardware. The experiment consists of a motor-propeller system powered via \gls{ESC}s, with \gls{PWM} signals generated by the ArduPilot flight controller installed on Pixhawk 6x hardware. Different \gls{PWM} levels were supplied from a 4-cell \gls{LiPo} battery, producing corresponding thrust and moment outputs. These forces were simultaneously measured using a force-torque sensor. Due to their differing mechanical characteristics, measurements were performed separately for the commercial rigid propellers and the Tombo propellers. The measurement results are summarized from the thrust and moment measurement experiment for the rigid and Tombo propellers.